\section{Discussion}
\label{sec:discussion}

\paragraph{Clustering is a poor labeler but a useful prior.}
Used directly, the 30-component mixture is the worst model in
Table~\ref{tab:p1models} (0.780 sample-F1, 0.570 macro-F1). Inspecting the
components explains why: they overwhelmingly separate Recruitment from
everything else, and additional components mostly subdivide Recruitment further
rather than isolating Safety or Efficacy. Since the embeddings do carry that
signal---a purely supervised model reaches Safety F1 0.947---the failure is the
clustering objective, not the features. Unsupervised structure in this space
tracks the dominant topic, and with Recruitment on 64\% of trials the dominant
topic is Recruitment. What does work is demoting the mixture from decision-maker
to prior: as a 0.2-weighted term inside the pseudo-labeling step it raises
cross-validated macro-F1 from 0.805 to 0.814 (and validation sample-F1 from 0.889
to 0.893) over identical self-training without it. The mixture's contribution is
to regularize which unlabeled trials get promoted, not to name their reasons.

\paragraph{Why the extra features fail.}
This is the ablation's most informative result. The reason embeddings are a
compressed encoding of the very stopping-reason text the labels were derived
from, so they are close to a sufficient statistic for the target. Criteria,
disease and drug fields describe \emph{what the trial was}, not \emph{why it
stopped}---a trial can be halted for funding reasons regardless of its indication.
Appending them adds 20--50 parameters per block to a model fit on 157 labeled
trials, so the variance cost arrives immediately while the bias reduction never
does: criteria TF-IDF drops validation sample-F1 from 0.893 to 0.876. The Morgan
fingerprints are the clearest null result, moving cross-validated macro-F1 by
$+0.001$---within repeat-to-repeat noise---which is unsurprising given that 44\%
of trials have no matched drug and receive an all-zero block. \texttt{start\_year}
fails for a more interesting reason: it is genuinely predictive in-sample, because
the labeled trials are deliberately the earliest ones, but the test trials are the
latest, so the model would be extrapolating a year effect fit entirely outside the
test range. Including everything is actively harmful here, not merely redundant.

\paragraph{Rare labels dominate macro-F1.}
Recruitment appears on 64\% of trials and Funding on 5\%. Funding F1 (0.400) and
Other F1 (0.500) rest on 4 and 10 positive validation trials respectively, so a
single trial moves macro-F1 by several points. This is why the validation macro-F1
column in Table~\ref{tab:p1models} is non-monotone with the CV column, and why we
selected on the latter: the final model has slightly lower validation macro-F1
than the balanced LR + $k$NN blend (0.714 vs.\ 0.724) but the best CV macro-F1
(0.814 vs.\ 0.800), and the 235-trial CV estimate is the more reliable of the two.

\paragraph{What the stack buys in Part 2.}
The fitted meta-model weights are informative: criteria LR 0.62, criteria $k$NN
0.55, drug-name LR 0.41, metadata 0.16, with disease (0.06) and ICD (0.03) nearly
ignored and drug-knowledge slightly negative ($-0.11$). The drug-knowledge block
is not useless in isolation---it is redundant once drug names are present, since
targets and pathways are a deterministic function of the matched drug, and the
meta-model subtracts it as a partial correction. The $k$NN's large weight shows
the stack is exploiting near-duplicate trials that the linear criteria model
cannot represent. Pooling phases matters most for Phase~I, which has only 1,012
training rows and gains the most from the shared representation.

\paragraph{The F1 / log-loss trade-off.}
The unweighted stack \emph{lowers} mean validation F1 from 0.739 to 0.736 while
improving log loss (0.591 $\to$ 0.559) and PR-AUC (0.796 $\to$ 0.810). This is not
a regression: F1 at a fixed 0.5 threshold is a statement about calibration near
that threshold, and the stack's meta-model shifts probabilities toward the base
rate, pushing borderline positives below 0.5. Ranking quality, which is what
ROC-AUC and PR-AUC measure, improves in every phase. Recency weighting then raises
the probabilities again and lifts F1 to 0.755, better than either predecessor.
Since the competition metric is log loss, we optimized for calibration and treat
the F1 recovery as a bonus; a project scored on F1 would instead have tuned the
threshold per phase.

\paragraph{Where the errors live.}
The strongest signal in Table~\ref{tab:p2errors} is that failure is hardest to
detect in Phase~III and among trials with no recorded stopping reason. Both follow
from the same cause: these are trials that ran to completion and missed their
endpoint, so nothing in the eligibility text, indication or drug marks them as
different from a successful trial---and Phase~III has a 0.65 success base rate
pulling predictions upward. By contrast, trials stopped for Recruitment or Funding
are administrative failures, and the stack catches them better (Phase~II
Recruitment 0.207 $\to$ 0.149) partly because small, slow-recruiting trials leave
traces in criteria length and drug-count metadata. The improvement is largest
exactly where the baseline was worst: Phase~I no-reason 0.643 $\to$ 0.321,
Phase~III no-reason 0.695 $\to$ 0.585.

\paragraph{How much of the error table to believe.}
Two caveats temper this. Category sizes are small---Safety, Funding, Efficacy and
Other have fewer than 25 failed trials in every phase, and Phase~III Efficacy has
8, so its apparent 1.000 $\to$ 0.875 $\to$ 1.000 trajectory is one trial moving.
Phase~III Funding has no trials at all. More subtly, 171 of the 198 matched reason
labels are Part~1 \emph{predictions} rather than ground truth, and the Part~1
model has roughly 10\% sample-F1 error. Those errors are not random: because the
Part~1 model over-predicts the majority reason, trials are more likely to be
misfiled \emph{into} Recruitment than out of it, which biases the Recruitment row
toward the population average and makes the rare-reason rows noisier than their
sample sizes alone suggest. Only the Phase~II Recruitment row and the three
no-reason rows (whose membership is determined by a join, not a prediction, and so
is error-free) support confident conclusions.

\paragraph{The validation-to-test gap.}
Mean log loss rises from 0.556 on validation to 0.58688 on Kaggle. The cause is
temporal: validation trials started up to 2013 and were split at random from the
training trials, while test trials started 2013--2022, later than every training
row, and the success rate climbs from 0.38 in 2000 to 0.81 in 2013. A model fit on
the pooled base rate is therefore systematically under-confident on test. Recency
weighting is our response and was the single largest holdout gain we found
(0.5514 $\to$ 0.5475 fitting to 2009; 0.4657 $\to$ 0.4553 fitting to 2011); it
redistributes mean test predictions per phase rather than inflating them
uniformly---Phase~III rises from 0.711 to 0.738 while Phase~I falls from 0.617 to
0.598 and Phase~II from 0.565 to 0.557---because recent Phase~III trials succeeded
far more often than recent early-phase ones. The residual gap reflects what we deliberately did not model. We
built an alternative submission that lets the meta-model learn an explicit
Phase~III year trend (Phase~III success rises 0.71 in 2009 to 0.96 in 2013), which
pushes mean Phase~III test predictions to 0.788; on the holdouts it gains only
about 0.001 log loss, and it extrapolates a trend fit on four years out to 2022,
so we kept it as a recorded alternative rather than the submission. A per-phase
trend for all phases was rejected outright: it fits flat or negative Phase~I/II
slopes and would drive those predictions below the training base rate. Beyond the
year effect, the test trials differ in disease and drug mix, and no amount of
reweighting on pre-2014 data can calibrate for that.
